Do sklářského kmene se přidávají:
\begin{itemize}
\item \textbf{Urychlovače tavení:} urychlují tavení či čeření; borax, kyselina boritá, fluoridy, křemičitan sodný.
\item \textbf{Čeřiva:} uvolňují plyny, které odnášejí bubliny ze skloviny; Na$_2$SO$_4$ nebo směs As$_2$O$_3$ a KNO$_3$.
\item \textbf{Barviva:} způsobují zbarvení rozpuštěnými ionty či molekulami nebo koloidními částicemi.
\item \textbf{Odbarviva:} zajišťují čirost a průzračnost; arsenik, ledek, MnO$_2$.
\item \textbf{Kaliva:} částice s odlišným indexem lomu způsobují zákal; např. kryolit Na$_3$AlF$_6$.
\end{itemize}
